class of @@M0@@ belongs to @@M1@@, the three coefficients @@M2@@ remain bounded by 1.13.

Let @@M3@@ be a multiple of @@M4@@ which is very ample. Then, in the Hilbert polynomials @@M5@@ relative to @@M6@@ of the @@M7@@ in @@M8@@, the three coefficients @@M9@@ remain bounded by 2.10. Consequently, @@M10@@ is bounded by 2.11\,; this gives the first part of 3.8.

As for the second part of 3.8, the ``generic'' representability of the Picard functors was established in XII 1.2, as was their separatedness. By an argument similar to that of XII 2, using 3.9 and XII 1.1, one reduces to the case of @@M11@@ normal over @@M12@@. Now, for the morphism @@M13@@ to be quasi-affine, it suffices that it be (separated and) quasi-finite (EGA IV 8.12.8), which can be checked on the geometric fibers, and thus follows from the following lemma:

\sgaStableTarget{sga6.e13.lem.3-11}\textbf{Lemma 3.11.} \emph{Let @@M14@@ be an algebraically closed field, @@M15@@ a normal projective @@M16@@-scheme endowed with an ample sheaf @@M17@@, and @@M18@@ the zero scheme of a section @@M19@@ of @@M20@@. We suppose that @@M21@@ has dimension @@M22@@ (resp. @@M23@@). Then the morphism @@M24@@ (resp. the morphism @@M25@@ (see 4.1)) is of finite type and has a kernel @@M26@@ which is a finite group (in general nonreduced).}

Indeed, one may evidently suppose @@M27@@\,; thus @@M28@@ is regular, @@M29@@ being integral. The morphism under consideration is of finite type by the first part of 3.8 (resp. by 4.7 (iii)). Consequently, the kernel @@M30@@ is an algebraic group, hence contained in @@M31@@, and is therefore the same
